\uuid{WHpC}
\exo7id{4049}
\titre{exo7 4049}
\auteur{quercia}
\organisation{exo7}
\datecreate{2010-03-11}
\isIndication{false}
\isCorrection{false}
\chapitre{Développement limité}
\sousChapitre{Equivalents}
\cpge{mpsi-mp2i.fonctions-valeurs-reelles-complexes.c, pcsi.techniques-calcul-fonctions.c, ptsi.techniques-calcul-fonctions.c, mpsi-mp2i.analyse-asymptotique.a, pcsi.analyse-asymptotique.a, ptsi.analyse-asymptotique.a, mpsi-mp2i.analyse-asymptotique.b, pcsi.analyse-asymptotique.b, ptsi.analyse-asymptotique.b, bcpst1.limites-continuite.b, bcpst1.developpements-limites.a, bcpst1.developpements-limites.b}
\module{Analyse}
\niveau{L1}
\difficulte{}

\contenu{
\question{
Soient $u,v$ deux fonctions positives, $u \sim v$, $u \to 0$.
Montrer que $\frac {u^v - v^u}{u-v} \sim -\ln(v)$.

(\'Ecrire $u = v + w$ avec $w/v \to 0$)
}
}
